\documentclass[10pt,a4paper]{amsart}
\usepackage[margin=19mm]{geometry}
\usepackage{amsmath,amssymb,mathtools}
\usepackage[scr=rsfs,cal=euler]{mathalfa}
\usepackage{xfrac}
\usepackage[unicode,hidelinks,bookmarks=false]{hyperref}
\newcommand{\ie}{i.e.\ }
\newcommand{\cf}{cf.\ }
\newcommand{\resp}{respectively\ }
\newcommand{\Subsection}{\subsection}
\providecommand{\nobreakdash}{\nobreak\hbox{-}}
\setlength{\emergencystretch}{2em}
\multlinegap0pt
\usepackage{bm}
\usepackage{bbm}
\usepackage{dsfont}
\usepackage{xspace}
\usepackage{cancel}
\usepackage{mathtools}
\usepackage{amsthm}
\makeatletter
\@ifundefined{fact}{\newtheorem{fact}{Fact}}{}
\@ifundefined{lemma}{\newtheorem{lemma}{Lemma}}{}
\makeatother
\newcommand{\E}{\mathbb{E}}
\newcommand{\N}{\mathbb{N}}
\newcommand{\visible}{\leftrightarrow{}}
\renewcommand{\P}{\mathbb{P}}
\title[Exact Recovery in the Geometric Hidden Community Model]{Exact Recovery in the Geometric Hidden Community Model}
\author{Julia Gaudio, Andrew Jin}
\date{Source: arXiv:2601.17591v1 (CC BY 4.0)}
\begin{document}
\maketitle

\noindent\emph{Source and licence.} Exact Recovery in the Geometric Hidden Community Model. Version arXiv:2601.17591v1 (CC BY 4.0), distributed under the Creative Commons Attribution 4.0 International licence, as recorded in the arXiv licence metadata for this exact version and shown on the abstract page https://arxiv.org/abs/2601.17591v1. Full rights notice: licenses/arxiv-2601.17591-CC-BY-4.0.txt.\par\smallskip

\noindent\textit{Editorial scope.} Counted unit: the Lemma whose source label is \texttt{lem:correct\_labels\_MAP\_failures} in the article (source0.tex lines 979--989, source label \texttt{lem:correct\_labels\_MAP\_failures}), delivered together with the article's own complete original proof of it (source0.tex lines 990--1064, one \texttt{proof} environment of 75 lines). No mathematics is written, repaired or re-derived: statement and proof are the article's own text line for line. Required context retained uncounted: fact:CGF\_formula (lines 972--975). Every definition, display and label of those lines is retained unchanged. Omitted from this package and not redistributed: 1--971, 976--978, 1065--1463. Nothing inside a retained excerpt is omitted or altered: every retained body is delivered exactly as the article's lines are written, so no omission receipt of any kind accompanies this package. Citations: the retained excerpts carry no citation command at all, so no bibliography entry of the article is transcribed and no reference is redistributed. Reference adaptation: none was needed and none was made; every reference inside the delivered unit resolves inside it, so no unresolved reference or citation is produced. Printed slips kept literal: no printed word, symbol, hypothesis or number of the counted statement or of its proof was altered. Layout and class: the article's own class is \texttt{article} with packages including graphicx, geometry, amsfonts, enumitem, amsmath, amssymb, amsthm, mathtools, comment, hyperref, xcolor, algorithm, algpseudocode, bbm; the wrapper is the lane's pinned \texttt{amsart} editorial wrapper at 10pt on A4, which restates the theorem-like environments the retained text needs and adds the article's own macro definitions that the retained text needs (\texttt{\textbackslash E}, \texttt{\textbackslash N}, \texttt{\textbackslash P}, \texttt{\textbackslash visible}), with extra wrapper packages (bm, bbm, dsfont, xspace, cancel, mathtools). The delivered numbering is the unit's own. Assets: no retained span contains a figure environment, a graphics-inclusion command, a PDF-page inclusion, a table or a TikZ picture; no image file is copied into this package, the package declares no asset dependency and no image file is redistributed with it. Licence: version arXiv:2601.17591v1 of this article is distributed under the Creative Commons Attribution 4.0 International licence, as recorded in the arXiv licence metadata for this exact version and shown on the abstract page https://arxiv.org/abs/2601.17591v1. A full rights notice is delivered at licenses/arxiv-2601.17591-CC-BY-4.0.txt. Characters: every retained excerpt is pure ASCII, so the delivered engine, XeLaTeX with its default Latin Modern text font, reports no missing character.
\begin{fact}
\label{fact:CGF_formula}
    Let $Y$ be a random variable, and $X_1, X_2, \ldots, X_Y$ be $Y$ i.i.d. copies of a random variable $X$ independent of $Y$. Let $S = \sum_{i=1}^Y X_i$. Then, $\Lambda_S(t) = \Lambda_Y(\Lambda_X(t))$,
\end{fact}

\begin{lemma}
\label{lem:correct_labels_MAP_failures}
    If $\lambda \nu_d r^d \min_{i \neq j} D_+(\theta_i \Vert \theta_j; \pi, g) > 1$, then for a fixed $0 < \epsilon \leq (\lambda \nu_d r^d \min_{i \neq j} D_+(\theta_i \Vert \theta_j; \pi, g) - 1 ) / 2$, we have that
    \begin{equation*}
        \P(\ell_i(v, \sigma^*) - \ell_j(v, \sigma^*) \leq \epsilon \log n \mid \sigma^*(v) = i) = n^{-(1 + \Omega(1))}
    \end{equation*}
    where
    \begin{equation*}
        \ell_i(v, \sigma) = \sum_{v \in V: u \visible v} \log(\overline{p}_{i, \sigma(u)}(x_{uv};\|u - v\|)).
    \end{equation*}
\end{lemma}

\begin{proof}
    Fix a vertex $v \in V$ with $\sigma^*(v) = i$. We construct a random variable $X$ such that 
    \begin{equation*}
        X \sim \ell_j(v, \sigma^*) - \ell_i(v, \sigma^*).
    \end{equation*}
    Observe that
    \begin{align*}
         \ell_j(v, \sigma^*) - \ell_i(v, \sigma^*) &= \sum_{u \in V: u \visible v} \log\bigg(\frac{p_{j, \sigma^*(v)}(x_{uv}; \|u-v\|)}{p_{i, \sigma^*(v)}(x_{uv}; \|u-v\|)}\bigg) \\
         &= \sum_{s \in Z} \bigg( \sum_{u \in V: u \visible v, \sigma^*(u) = s} \log\Big(\frac{p_{js}(x_{uv}; \|u-v\|)}{p_{is}(x_{uv}; \|u-v\|)}\Big) \bigg)
    \end{align*}
    Now, we construct $X_s$ for each $s \in Z$ such that
    \begin{equation*}
        X_s \sim \log\bigg(\frac{p_{js}(x_{uv}; \|u-v\|)}{p_{is}(x_{uv}; \|u-v\|)}\bigg).
    \end{equation*}
    Let $D$ be a random variable with density
    \begin{equation*}
        f_D(x) :=
        \begin{cases}
            dx^{d-1}/r^d \text{ if } 0 \leq x \leq r \\
            0 \text{ otherwise,} 
        \end{cases}
    \end{equation*}
    which represents the distance between $v$ and a vertex $u$ placed uniformly at random in the neighborhood of $v$. Let $X_{is}$ be a random variable coupled to $D$, such that conditioned on $D = y$, we sample $X_{is}$ from $P_{is}(y)$. Note that $X_{is}$ represents the edge weight between $u$ and $v$. Then, we define
    \begin{equation*}
        X_s := \log\bigg( \frac{p_{js}(X_{is} ;D)}{p_{is}(X_{is} ;D)} \bigg).
    \end{equation*}  
    Now, letting $N_s \sim \text{Pois}(\lambda \nu_d r^d \pi_s \log n)$ and $\{X_s^{(m)}\}_{m \in \N} \overset{\text{iid}}{\sim} X_s$, we can define
    \begin{equation*}
        X = \sum_{s \in Z}\bigg( \sum_{m=1}^{N_s} X_s \bigg).
    \end{equation*}
    By the Chernoff bound, we have that for any fixed $\epsilon > 0$,
    \begin{align}
        \P(\ell_i(v, \sigma^*) - \ell_j(v, \sigma^*) \leq \epsilon \log n \mid \sigma^*(v) = i) &= \P(X \geq -\epsilon \log n) \nonumber \\
        &\leq \inf_{t \in [0, 1]} \E[e^{tX}] e^{t\epsilon \log n}  \nonumber \\
        &= \inf_{t \in [0, 1]} \E[e^{tX}]n^{t\epsilon}.     \label{eq:Chernoff_bound_X}
    \end{align}
    Thus, we need to compute the moment-generating function (MGF) of $X$. Observe that
    \begin{align}
    \label{eq:MGF_X}
        \E[e^{tX}] &= \E\bigg[\exp\bigg(t\sum_{s \in Z} \sum_{m=1}^{N_s} X_s\bigg)\bigg] \notag \\
        &= \prod_{s \in Z} \E\bigg[\exp\bigg(t\sum_{m=1}^{N_s} X_s\bigg)\bigg] \notag \\
        &= \prod_{s \in Z} \exp\Big(\Lambda_{N_s}(\Lambda_{X_s}(t))\Big).
    \end{align}
    Computing the MGF of $X_s$ using the law of total expectation, we obtain that
    \begin{align*}
        \E[e^{tX_s}] &= \E\bigg[\exp\bigg(t \log\Big( \frac{p_{js}(X_{is};D)}{p_{is}(X_{is};D)} \Big)\bigg)\bigg] \\
        &= \E\bigg[\E\bigg[\exp\bigg(t \log\Big( \frac{p_{js}(X_{is};D)}{p_{is}(X_{is};D)} \Big)\bigg) \mid D = y\bigg]\bigg] \\
        &= \int_0^r \E\bigg[ \exp\bigg(t \log\Big( \frac{p_{js}(X_{is};y)}{p_{is}(X_{is};y)} \Big)\bigg)\bigg] \frac{dy^{d-1}}{r^d} dy \\
        &= \int_0^r \phi_t(P_{js}(y), P_{is}(y)) \frac{dy^{d-1}}{r^d} dy \\
        &= \overline{\phi}_t(P_{js}, P_{is}).
    \end{align*}
    Hence, $\Lambda_{X_s}(t) = \log \overline{\phi}_t(P_{js}, P_{is})$. Now, combining the fact that $\Lambda_{N_s}(t) = (\lambda \nu_d r^d \pi_s \log n)(e^t - 1)$ with Fact \ref{fact:CGF_formula}, we obtain that
    \begin{align}
    \label{eq:CGF_calculation}
        \Lambda_{N_s}(\Lambda_{X_s}(t)) &= (\lambda \nu_d r^d \pi_s \log n) (\overline{\phi}_t(P_{js}, P_{is}) - 1) \notag \\
        &= - \lambda \nu_d r^d \pi_s (1 - \overline{\phi}_t(P_{js}, P_{is})) \log n.
    \end{align}
    Substituting \eqref{eq:CGF_calculation} into \eqref{eq:MGF_X} yields the MGF of $X$, which is
    \begin{align}
    \label{eq:MGF_X_finished}
        \E[e^{tX}] &= \prod_{s \in Z} \exp\bigg( - \lambda \nu_d r^d \pi_s (1 - \overline{\phi}_t(P_{js}, P_{is})) \log n \bigg) \notag \\
        &= \prod_{s \in Z} n^{- \lambda \nu_d r^d \pi_s (1 - \overline{\phi}_t(P_{js}, P_{is}))} \notag \\
        &= n^{-\lambda \nu_d r^d \big( 1 - \sum_{s \in Z} \pi_s \overline{\phi}_t(P_{js}, P_{is}) \big)}.
    \end{align}
    Finally, substituting \eqref{eq:MGF_X_finished} into \eqref{eq:Chernoff_bound_X}, we obtain that for any fixed $\epsilon > 0$,
    \begin{align*}
        \P(\ell_i(v, \sigma^*) - \ell_j(v, \sigma^*) \leq \epsilon \log n \mid \sigma^*(v) = i) &\leq \inf_{t \in [0, 1]} n^{-\lambda \nu_d r^d \big( 1 - \sum_{s \in Z} \pi_s \overline{\phi}_t(P_{js}, P_{is}) \big) + t\epsilon} \\
        &\leq n^{-\lambda \nu_d r^d \big( 1 - \inf_{t \in [0, 1]} \sum_{s \in Z} \pi_s \overline{\phi}_t(P_{js}, P_{is}) \big) + \epsilon} \\
        &= n^{-\big(\lambda \nu_d r^d D_+(\theta_i \Vert \theta_j; \pi, g) - \epsilon \big)}.
    \end{align*}
    Since $\lambda \nu_d r^d \min_{i \neq j} D_+(\theta_i \Vert \theta_j; \pi, g) > 1$ by the achievability criterion, we can take $\epsilon = (\lambda \nu_d r^d \min_{i \neq j} D_+(\theta_i \Vert \theta_j; \pi, g) - 1)/2$ to obtain that
    \begin{equation*}
        \P(\ell_i(v, \sigma^*) - \ell_j(v, \sigma^*) \leq \epsilon \log n \mid \sigma^*(v) = i) = n^{-(1 + \Omega(1))}.
    \end{equation*}
\end{proof}

\end{document}
